% topic5.tex -- Module 8 Topic 5: Deep inside the atom, Q57-70.
\topictitle{Topic 5 - Deep Inside the Atom}

% ---- Q57 (not fundamental)
\begin{quizq}{}
According to the Standard Model, which particle is NOT a fundamental particle?
\opts{Electron}{Muon}{Photon}{Proton}
\end{quizq}

% ---- Q58 (fundamental)
\Needspace{12\baselineskip}
\begin{quizq}{}
According to the Standard Model, which particle IS a fundamental particle?
\opts{Proton}{Tau}{Pion}{Neutron}
\end{quizq}

% ---- Q59 (quarks in a neutron)
\Needspace{12\baselineskip}
\begin{quizq}{}
Which combination of quarks makes up a neutron?
\opts{udd}{uud}{uuu}{ddd}
\end{quizq}

% ---- Q60 (charge of the omega-minus)
\Needspace{12\baselineskip}
\begin{quizq}{}
The omega-minus particle is made of three strange quarks. Each strange quark has a charge of
$-\frac{1}{3}e$. What is the charge of the omega-minus particle?
\opts{$-e$}{$-\frac{1}{3}e$}{$0$}{$+e$}
\end{quizq}

% ---- Q61 (beta decay at the quark level)
\Needspace{15\baselineskip}
\begin{quizq}{}
In beta-minus decay, a neutron (udd) changes into a proton (uud). Which particles are emitted?
\opts{A positron and an electron neutrino}{An electron and an electron antineutrino}
  {A gluon only}{A gamma-ray photon only}
\end{quizq}

% ---- Q62 (meson)
\Needspace{12\baselineskip}
\begin{quizq}{}
What is a meson made of?
\opts{Three quarks}{Three antiquarks}{Two quarks}{A quark and an antiquark}
\end{quizq}

% ---- Q63 (force carriers)
\Needspace{20\baselineskip}
\begin{quizq}{}
Which row correctly shows the exchange particle (force carrier) for the strong force and for the
electromagnetic force?
\begin{center}
  \renewcommand\arraystretch{1.7}
  \begin{tabular}{@{}p{10mm}|>{\centering\arraybackslash}p{45mm}|>{\centering\arraybackslash}p{45mm}|}
    \cline{2-3}
    & \textit{Strong force} & \textit{Electromagnetic force}\\ \cline{2-3}
    A. & Photon & Gluon\\ \cline{2-3}
    B. & W and Z bosons & Photon\\ \cline{2-3}
    C. & Gluon & W and Z bosons\\ \cline{2-3}
    D. & Gluon & Photon\\ \cline{2-3}
  \end{tabular}
\end{center}
\end{quizq}

% ---- Q64 (electron-positron collision making a proton-antiproton pair)
\Needspace{18\baselineskip}
\begin{quizq}{}
Can a collision between an electron and a positron produce a proton and an antiproton?
\opts{No, because electrons and positrons are leptons, so they cannot produce other kinds of particle.}
  {No, because mass cannot be created in a collision.}
  {Yes, if the total energy is at least $2m_\mathrm{p}c^2$, because charge is conserved and a particle is
   made with its antiparticle.}
  {Yes, at any energy, because charge is conserved.}
\end{quizq}

% ---- Q65 (minimum energy for a proton-antiproton pair)
\Needspace{12\baselineskip}
\begin{quizq}{}
What is the minimum energy needed to create a proton--antiproton pair?
\opts{\qty{1.0e-18}{J}}{\qty{1.5e-10}{J}}{\qty{3.0e-10}{J}}{\qty{6.0e-10}{J}}
\end{quizq}

% ---- Q66 (why accelerators need high energies)
\Needspace{18\baselineskip}
\begin{quizq}{}
Why do particle accelerators need to give particles very high energies to discover new particles?
\opts{So that the particles travel faster than light.}
  {The energy of a collision can be converted into the mass of new, heavy particles ($E = mc^2$), and
   high-energy particles have short de Broglie wavelengths, so they can probe very small distances.}
  {High energies increase the charge on the particles.}
  {So that the particles can escape Earth's gravity.}
\end{quizq}

% ---- Q67 (conservation of charge)
\Needspace{14\baselineskip}
\begin{quizq}{}
Which reaction is NOT possible, because it does not conserve electric charge?
\opts{$\mathrm{p} + \mathrm{p} \rightarrow \mathrm{p} + \mathrm{p} + \pi^0$}
  {$\mathrm{p} + \mathrm{n} \rightarrow \mathrm{p} + \mathrm{p} + \pi^-$}
  {$\mathrm{p} + \mathrm{p} \rightarrow \mathrm{p} + \mathrm{n} + \pi^+$}
  {$\mathrm{n} \rightarrow \mathrm{p} + \mathrm{e}^+ + \nu_\mathrm{e}$}
\end{quizq}

% ---- Q68 (evidence that protons are not fundamental)
\Needspace{16\baselineskip}
\begin{quizq}{}
What evidence showed that protons are not fundamental particles?
\opts{High-energy electrons scattered off protons at large angles, showing that protons contain
   smaller, point-like charged particles.}
  {Protons are deflected by magnetic fields.}{Protons have a positive charge.}
  {Protons are found in every nucleus.}
\end{quizq}

% ---- Q69 (no strong force)
\Needspace{12\baselineskip}
\begin{quizq}{}
Which particle does NOT feel the strong nuclear force?
\opts{Proton}{Electron}{Neutron}{Pion}
\end{quizq}

% ---- Q70 (annihilation photons)
\Needspace{13\baselineskip}
\begin{quizq}{}
An electron and a positron, both almost at rest, annihilate each other and produce two gamma-ray
photons. What is the energy of each photon?
\opts{\qty{2.7e-22}{J}}{\qty{4.1e-14}{J}}{\qty{8.2e-14}{J}}{\qty{1.6e-13}{J}}
\end{quizq}
